\paragraph{About this part.} Every earlier part ends with the same question: \emph{is the model any
good?} Part~VI answers it. Evaluation in this field fails in predictable ways---a metric that rewards
the wrong thing, a test set that leaks information from the training set, or a winner chosen from
thousands of tries that only looks good by chance---and each section addresses one of them. In order:
\begin{itemize}[leftmargin=1.5em,itemsep=1pt]
\item which \emph{statistical} scores measure forecast quality honestly, and why accuracy and F1 are
  not enough (Section~\ref{sec:stat-metrics});
\item how to measure \emph{economic} value---profit after costs, fills and adverse selection---which is
  what a forecast is ultimately for (Section~\ref{sec:econ-metrics});
\item how to split time-ordered data so that the test really is unseen (Section~\ref{sec:validation});
\item why more complex models need more evidence, and how to correct for having tried many models
  (Section~\ref{sec:generalisation});
\item how to choose between models of different complexity, which turns Occam's razor (Principle~3)
  into a procedure (Section~\ref{sec:parsimony});
\item and what a benchmark result must report for others to compare and reproduce it
  (Section~\ref{sec:benchmark}).
\end{itemize}
The thread through all six is the distinction of Part~I: a forecast is judged first as a forecast, then
as a trading input, and the second judgement is the one that matters.

% ================================================================
\section{Statistical metrics}
\label{sec:stat-metrics}
% ================================================================

The right statistical metric depends on the target (Table~\ref{tab:metrics}). These metrics
measure predictive quality; none of them, by itself, measures trading value.

\begin{table}[!htbp]
\centering\small
\caption{Statistical metrics by prediction target.}
\label{tab:metrics}
\begin{tabularx}{\linewidth}{lX}
\toprule
Target & Metrics \\
\midrule
Direction & log loss, Brier score~\citep{brier1950}, AUC~\citep{fawcett2006}, balanced accuracy, calibration error~\citep{guo2017} \\
Direction with confidence & risk--coverage curves; F1 by confidence tier; interval coverage against
  nominal (e.g.\ the 68\% coverage reported by~\citet{manoharan2026uqlob}) \\
Return & MSE, MAE, CRPS~\citep{gneiting2007} \\
Event time or fill & right-censored log-likelihood~\citep{arroyo2024}; calibration by queue-position bucket; concordance with caution \\
Intensity & point-process log-likelihood; time-rescaling residual tests \\
\bottomrule
\end{tabularx}
\end{table}

\paragraph{Proper scoring rules.} A scoring rule is \emph{proper} if a forecaster minimises the expected
score by reporting their true beliefs~\citep{gneiting2007}. Log loss, the Brier score and the CRPS are
proper; accuracy and F1 are not, because they reward only the side of a threshold a forecast falls on.
For any model whose output will be used as a probability---for sizing, for the confidence gating of
Section~\ref{sec:tradeability}, or for the fill-probability term of execution value---a proper score
should be the primary statistical metric, and calibration should be checked directly (events forecast
at 70\% should happen about 70\% of the time). For fill-time models, censoring by cancellations makes
some familiar survival metrics improper; the right-censored log-likelihood is the safe
choice~\citep{arroyo2024}.

% ================================================================
\section{Economic metrics}
\label{sec:econ-metrics}
% ================================================================

Economic evaluation starts from the aggressive and passive returns of Section~\ref{sec:passive} and
the decomposition 
\[
  \E[\Pi^{\text{pass}}] = \P(\Fill)\,\E[\Pi^{\text{pass}} \mid \Fill],
\]
which separates execution probability
from conditional execution quality. For passive strategies both factors should be reported, broken
down by queue position where the data allow.

Table~\ref{tab:metric-levels} orders evaluation metrics by how close they come to trading value. A
study should say which levels it reaches; most LOB papers stop at the first or second.

\begin{table}[!htbp]
\centering\footnotesize
\caption{Levels of evaluation, from statistical to economic.}
\label{tab:metric-levels}
\begin{tabularx}{\linewidth}{l>{\raggedright\arraybackslash}X>{\raggedright\arraybackslash}X}
\toprule
Level & Question & Typical metrics \\
\midrule
Prediction & Is the direction or value right? & accuracy, F1, AUC, MSE \\
Probabilistic prediction & Are the stated probabilities right? & log loss, Brier, CRPS, calibration \\
Execution & Will the order trade? & fill probability, fill-time likelihood \\
Conditional execution & How good is the trade when it happens? & mark-outs, adverse selection by queue position \\
Trading & Does it make money after costs? & net P\&L, Sharpe ratio, transaction-level overlap $\mathrm{TO}$~\citep{briola2024} \\
Capacity & How much can be traded before impact erases it? & turnover, impact-adjusted P\&L \\
Robustness & Does it survive new periods, assets and regimes? & out-of-sample by regime and asset \\
Statistical validity & Could the result be luck given the search? & deflated Sharpe ratio, null maximum, backtest-overfitting probability \\
\bottomrule
\end{tabularx}
\end{table}

Cost-aware directional metrics, which ask whether a predicted trade would have produced a favourable
executable outcome rather than whether the mid moved in the predicted direction, are an improvement
over raw accuracy. They are not a complete economic metric: they omit market impact, capacity,
inventory, repeated trading, the distribution of latency, and queue position in detail. In
Example~\ref{ex:as}, a metric that ignored $\P(\Fill)$ would misjudge the two positions. Such
metrics should be reported alongside realised trading metrics, not instead of them.

% ================================================================
\section{Temporal validation}
\label{sec:validation}
% ================================================================

Cross-validation with randomly assigned folds is inappropriate for most LOB problems because observations are
dependent in time. A chronological split preserves the direction of information flow, and
walk-forward evaluation repeats it. Where features and labels overlap in time, training samples
whose information overlaps the test period must be purged and an embargo imposed after each test
window~\citep{lopezdeprado2018}. The embargo should follow from the maximum look-back and the
forecast horizon, not be chosen arbitrarily.

\begin{example}[Sizing an embargo]
Features use a 60-second look-back and labels a 10-second horizon. The embargo at each
train/test boundary is at least 70 seconds. A walk-forward schedule might train on 20 sessions,
validate on 5 and test on 5, then roll forward by 5.
\end{example}

When several models are compared on the same data, a \emph{model confidence set}~\citep{hansen2011mcs}
identifies the subset that cannot be statistically distinguished from the best, at a stated
confidence level; it is the honest alternative to declaring the model with the highest score the
winner. Beyond a single test period, models should be stress-tested across volatility regimes, tick
and fee changes, and assets of different tick size, since predictability differs systematically
between large- and small-tick instruments~\citep{briola2024,gould2016qi}.

LOB research has large implicit search spaces: feature definitions, look-backs, horizons,
normalisations, thresholds, architectures, hyper-parameters, assets, sessions and execution
assumptions. When thousands of configurations are tried, the best one looks good by chance. Model
complexity and research-process complexity must be considered together
(Section~\ref{sec:search}).

% ================================================================
\section{Capacity, generalisation and multiple testing}
\label{sec:generalisation}
% ================================================================

A model that fits the training data well can still fail on new data, and the more flexible the
model, or the more models were tried, the more likely that is. This section gives the theory behind
that statement in three steps: why simpler model families generalise more reliably, how to correct a
performance measure for the number of models tried, and why simple models are also easier to test
honestly.

\subsection{Generalisation of simpler models}

Statistical learning theory makes the intuition behind Occam's razor precise. For a classifier $M$ chosen
from a family of models with \emph{VC dimension} $\mathrm{VC}$---roughly, the largest number of points
the family can label in every possible way---and trained on $N_s$ independent examples, with probability
at least $1-\varsigma$ the error rate on new data satisfies~\citep{vapnik1998}
\begin{equation}
  \text{err}(M) \le \widehat{\text{err}}_{N_s}(M) +
  \sqrt{\frac{\mathrm{VC}\,\big(\ln(2N_s/\mathrm{VC}) + 1\big) - \ln(\varsigma/4)}{N_s}} .
\end{equation}
\inwords{error on new data is at most the training error plus a penalty that grows with the
richness of the model family and shrinks with the amount of data. A richer family can fit more
patterns, including patterns that are only noise, and needs more data before its training error can be
trusted.} Rademacher complexity~\citep{bartlett2002}, which measures how well a model family can fit
random noise on the data actually observed, gives bounds of the same form that are never much worse than the VC bounds and can be tighter.

Two features of LOB data make the penalty bite. The signal is weak relative to the noise, so a family
that can fit noise will. And observations are strongly dependent in time, so the \emph{effective}
number of independent examples is far smaller than the number of events: a day of ES data has
millions of messages but far fewer independent price moves.

Two caveats keep the argument honest. Modern deep networks can fit even random labels perfectly and
still generalise on real labels~\citep{zhang2017rethinking}, and test error can fall again as capacity
grows past the point where the training data are fitted exactly (``double
descent'')~\citep{belkin2019}. Classical capacity bounds therefore do not decide by themselves whether
a deep model will generalise on LOB data. What decides it is out-of-sample evidence---which is why the
benchmark studies of Section~\ref{sec:benchmarks-dl} carry more weight than any bound.

\paragraph{Practical penalties.} Information criteria trade fit against parameter count explicitly:
the Akaike criterion adds twice the number of parameters to minus twice the log-likelihood~\citep{akaike1974},
and the Bayesian criterion adds the number of parameters times the logarithm of the sample
size~\citep{schwarz1978}, penalising complexity more heavily in large samples. The lasso adds a penalty
on the sum of absolute coefficients, which drives unhelpful coefficients exactly to zero and so selects
features automatically~\citep{tibshirani1996}. All three are applied to a single fit; none accounts for
the search that produced the model family in the first place, which is the subject of the next
subsection.

\subsection{Multiple testing and the deflated Sharpe ratio}

If $N_{\text{trials}}$ strategies with no skill are tested, and their estimated Sharpe ratios have
standard deviation $\sigma_{\text{SR}}$ across trials, the expected best estimate is
approximately~\citep{bailey2014}
\begin{equation}
  \E\big[\max \widehat{\text{SR}}\big] \approx \sigma_{\text{SR}}
  \left[(1-\gamma_E)\,\Phi^{-1}\!\left(1 - \frac{1}{N_{\text{trials}}}\right)
  + \gamma_E\, \Phi^{-1}\!\left(1 - \frac{1}{N_{\text{trials}}\, \mathrm{e}}\right)\right],
\end{equation}
where $\Phi^{-1}$ is the inverse of the standard normal distribution function, $\gamma_E \approx 0.5772$
is the Euler--Mascheroni constant and $\mathrm{e}$ is Euler's number.
\inwords{this is the Sharpe ratio that the luckiest of $N_{\text{trials}}$ worthless strategies is
expected to show. It grows with the number of trials, roughly like $\sqrt{2\ln N_{\text{trials}}}$; the
example of Section~\ref{sec:search} used that rough form.} Write this expected maximum as
$\text{SR}_0$. The \emph{deflated Sharpe ratio} is the probability that the selected strategy's true
Sharpe ratio exceeds $\text{SR}_0$, given its estimate $\widehat{\text{SR}}$ from $N_{\text{obs}}$
returns:
\begin{equation}
  \text{DSR} = \Phi\!\left(\frac{(\widehat{\text{SR}} - \text{SR}_0)\sqrt{N_{\text{obs}} - 1}}
  {\sqrt{1 - \widehat{\text{skew}}\;\widehat{\text{SR}} + \frac{\widehat{\text{kurt}} - 1}{4}\,\widehat{\text{SR}}^{\,2}}}\right),
\end{equation}
where $\widehat{\text{skew}}$ and $\widehat{\text{kurt}}$ are the sample skewness and (raw, not excess)
kurtosis of the returns and the Sharpe ratios are per period, not annualised.
\inwords{the numerator asks how far the observed Sharpe ratio is above what luck alone would produce
after this many trials; the denominator widens the uncertainty when returns are skewed or
fat-tailed. A DSR near 1 means the result is unlikely to be luck; near 0.5 or below, it is
indistinguishable from the best of a batch of noise.}

The same authors turn this into a budget. For a given length of history there is a maximum number of
configurations that can be tried before a strategy with no skill is expected to show an impressive
in-sample Sharpe ratio; with five years of data, ``no more than forty-five independent model
configurations should be tried'' if one wants to avoid an expected in-sample Sharpe ratio of 1 from a
strategy whose out-of-sample expectation is zero~\citep{bailey2014pseudo}. The probability of backtest
overfitting~\citep{bailey2014pbo}, White's reality check~\citep{white2000}, raised significance
thresholds~\citep{harvey2016} and false-discovery-rate control~\citep{benjamini1995} address the same
problem from different directions.

\subsection{Backtestability}

These tools favour simple models for a reason beyond generalisation: a simple model family is
\emph{backtestable}. When the family is small, the number of effective trials can be counted or
bounded, the null maximum can be computed, and the path from signal to decision can be audited. As the
family grows---deep networks with many hyper-parameters, richly parameterised point processes, or
automated search over model structures---the effective number of trials becomes hard to state, and with
it the bar a result must clear. Parsimony is therefore a property of the research process as well as of
the model.

% ================================================================
\section{Occam's razor as a model-selection principle}
\label{sec:parsimony}
% ================================================================

This section makes Principle~3 operational. We do not propose a single scalar penalty for
complexity. With $V(M)$ as in
Section~\ref{sec:notation}, let
\begin{equation}
  \mathcal{C}(M) = \left(\mathcal{C}_{\text{param}}(M),\, \mathcal{C}_{\text{lat}}(M),\,
  \mathcal{C}_{\text{fit}}(M),\, \mathcal{C}_{\text{mon}}(M),\, N_{\text{trials}}(M)\right)
\end{equation}
collect the dimensions of complexity: the number of fitted parameters, the inference latency, the
effort to fit and refit its parameters, the effort to monitor in production, and the number of
candidates $N_{\text{trials}}(M)$ searched to find the model (Section~\ref{sec:lamd}).
\inwords{complexity is a list of separate costs, not one number: how many knobs the model has, how
slow it is, how much work it is to tune and to watch, and how hard someone searched to find it.} A complex model
$M_1$ should be preferred to a baseline $M_0$ only if $\Delta V = V(M_1) - V(M_0)$ is statistically
credible and economically meaningful relative to 
\[
  \Delta\mathcal{C} = \mathcal{C}(M_1) - \mathcal{C}(M_0).
\]
Because the
components of $\mathcal{C}$ share no common unit, the framework asks for them to be reported separately
(Table~\ref{tab:parsimony}), so the reader can judge whether the added complexity is justified.

\begin{table}[!htbp]
\centering\small
\caption{Complexity dimensions of the model classes discussed in this survey, in orders of
magnitude. A comparison should add, for each row, the out-of-sample Sharpe ratio, the deflated
Sharpe ratio and the number of configurations tried, all under the same fill model.}
\label{tab:parsimony}
\begin{tabular}{llll}
\toprule
Model & Parameters & Typical inference latency & Fitting burden \\
\midrule
Model-free benchmark (\S\ref{sec:shadow}) & $\sim 10$ settings & $\mu$s & none \\
OFI baseline & $\sim 10$ & ns--$\mu$s & low \\
Shallow network & $10^3$--$10^4$ & $\mu$s & medium \\
Transformer & $\sim 10^6$ & ms & high \\
Hybrid & $\sim 10^6$ & ms & high \\
\bottomrule
\end{tabular}
\end{table}

\paragraph{An explicit exchange rate.} When a decision must be made, the separate dimensions can be
combined by stating, in advance, how much economic value each unit of added complexity must buy.
Writing $\Delta\mathcal{C}$ as a weighted sum of its components with weights chosen by the user, the
rule ``adopt $M_1$ only if $\Delta V > \chi\,\Delta\mathcal{C}$'' for a stated $\chi > 0$ makes
the trade-off explicit and auditable. The value of the rule lies in stating the weights before the
comparison, not in any particular choice of them.

\paragraph{Starting points ordered by complexity.} Table~\ref{tab:selection} orders candidate approaches by increasing
complexity, with the condition under which each is the natural starting point. It is a guide to where
the burden of proof lies, not a ranking of methods: a more complex row is justified when the simpler one
above it has been tried, on the same data and fill model, and found wanting.

\begin{table}[!htbp]
\centering\footnotesize
\caption{Starting points ordered by increasing complexity.}
\label{tab:selection}
\begin{tabularx}{\linewidth}{>{\raggedright\arraybackslash}p{0.24\linewidth}>{\raggedright\arraybackslash}X>{\raggedright\arraybackslash}X}
\toprule
Situation & Start with & Why \\
\midrule
Decision needed within microseconds; large-tick instrument & linear OFI or queue imbalance; queue-race probability (Section~\ref{sec:cdl}) & few operations, deterministic latency, strongest evidence in large-tick names~\citep{gould2016qi,briola2024} \\
Passive execution; fill probability matters & residual-volume and OFI baseline; then a survival model with queue position as an input~\citep{kanformer} & prices fill probability and adverse selection directly \\
Decision within milliseconds; liquid instrument & neural network on order-flow inputs~\citep{kolm2023,lucchese2024} & order-flow inputs carry most of the gain; in re-tests the simple MLP degraded least~\citep{prata2024} \\
Several horizons at once & sequence model only if baselines fail deflated and purged tests & complexity must clear the multiple-testing bar \\
Unusual volatility or liquidity & volatility gating of a simpler model & simple and robust when learned relations break \\
Several related instruments & lagged cross-asset OFI~\citep{contcucuringu2023} & modest complexity; evidence of forecasting value \\
\bottomrule
\end{tabularx}
\end{table}

\paragraph{A parsimony checklist.} The contrast between an over-engineered and a parsimonious approach
can be stated dimension by dimension (Table~\ref{tab:checklist}). The right-hand column is a default,
to be departed from with evidence.

\begin{table}[!htbp]
\centering\footnotesize
\caption{Over-engineered versus parsimonious defaults.}
\label{tab:checklist}
\begin{tabularx}{\linewidth}{l>{\raggedright\arraybackslash}X>{\raggedright\arraybackslash}X}
\toprule
Dimension & Over-engineered & Parsimonious default \\
\midrule
Features & many correlated indicators & a few interpretable drivers (OFI, imbalance, queue position) \\
Model & deep network chosen first & linear or logistic model first, then a shallow network \\
Fitting criterion & in-sample fit & out-of-sample economic value, with information criteria or shrinkage \\
Search & unrecorded trials & counted trials and a deflated result \\
Regime behaviour & fragile under shocks & degrades gracefully; gated by volatility \\
Diagnosis & opaque when it fails & failures traceable to an identifiable input \\
\bottomrule
\end{tabularx}
\end{table}

% ================================================================
\section{Benchmarking standard and datasets}
\label{sec:benchmark}
% ================================================================

A scientifically useful LOB benchmark should report at least the following.
\begin{description}[leftmargin=1.5em,style=nextline]
\item[Data] exchange; instrument; date range; number of sessions; L1, L2 or L3; event count; tick
  size; spread; depth.
\item[Prediction] target; horizon; sampling clock; feature look-back; normalisation.
\item[Model] architecture; parameter count; training procedure; hyper-parameter search;
  computational budget.
\item[Validation] chronological split; walk-forward procedure; purging; embargo; number of trials.
\item[Execution] latency; fees; spread; queue position; fill model; market-impact assumption.
\item[Results] statistical metric; economic metric; confidence interval; null benchmark;
  robustness across regimes.
\end{description}

Public datasets and frameworks each have limits (Table~\ref{tab:datasets}). No single benchmark
represents ``the LOB''; the strongest evidence survives across datasets, assets and regimes.

\begin{table}[!htbp]
\centering\footnotesize
\caption{Data sets and frameworks used in the literature surveyed.}
\label{tab:datasets}
\begin{tabularx}{\linewidth}{>{\raggedright\arraybackslash}p{0.22\linewidth}>{\raggedright\arraybackslash}X>{\raggedright\arraybackslash}X}
\toprule
Resource & What it provides & Limits \\
\midrule
FI-2010~\citep{ntakaris2018} & 5 Helsinki stocks, 10 days of 2010, 10 levels, normalised and downsampled & short, one market, pre-processed; results do not transfer reliably~\citep{prata2024} \\
LOBSTER~\citep{lobster} & Nasdaq books reconstructed from ITCH messages, any level, any period & commercial; Nasdaq only; one venue of a fragmented market \\
Exchange MBO feeds (Nasdaq ITCH~\citep{itch50}, CME MDP~3.0~\citep{cme_mdp3}) & full message-level data, queue positions & require capture or purchase; large; hard to share \\
Cryptocurrency exchange APIs & free, 24/7, many venues & different fees, participants and hours; often snapshots plus updates \\
LOBCAST~\citep{prata2024} & re-implementations of 15 models, common protocol & classification focus \\
LOBFrame~\citep{briola2024} & open framework with transaction-level evaluation & DeepLOB-centred \\
LOB-Bench~\citep{lobbench2025} & evaluation suite for generative LOB models & generation, not trading \\
\bottomrule
\end{tabularx}
\end{table}

A queue-exact replay can supply fields that a snapshot backtest cannot---queue position at entry,
the fill rule, the three latencies, the impact assumption---and in doing so makes the
\emph{Execution} block of the checklist concrete rather than nominal.
